\section{Design of \method}\label{sec:design}
\noindent\textbf{Overview.}
\method organizes functionality around \emph{risks} as the primary abstraction.
For each risk, we implement three pluggable module types:
\emph{(i)} attacks which realize the risk,
\emph{(ii)} defenses which protect against the risk, and
\emph{(iii)} metrics which evaluate susceptibility to the risk.
This modular structure is illustrated in Figure~\ref{fig:overview} and ensures that new components can be added consistently across domains.

\noindent\textbf{Intended users.}
\method is designed to serve two groups:
\begin{enumerate*}[label=\roman*),itemjoin={,\xspace}] % chktex 9 % chktex 10
    \item \emph{researchers}, who investigate new interactions, analyze how attack, defense, and metric variants affect them, and design defenses to mitigate or leverage these interactions
    \item \emph{practitioners}, who analyze trade-offs in deploying ML defenses specific to their domain and regulations.
\end{enumerate*}

\noindent\textbf{Scope.}
This design directly addresses the desiderata in Section~\ref{sec:problem}.
Below, we demonstrate how \method satisfies D1--D2 (comprehensive, consistent) by construction.
We evaluate D3 (extensibility) and D4 (applicability) empirically in Sections~\ref{sec:extensibility} and~\ref{sec:evaluation}.

\input{tables/tab_modules}

\subsection{Comprehensive}

\method spans eight risks across security, privacy, and fairness, with representative attacks, defenses, and metrics summarized in Table~\ref{tab:modules}.
We deliberately favor breadth by implementing widely used and generalizable modules across all risks, rather than depth in any single risk category.
This approach ensures coverage of diverse risk types while leaving room for the open-source community to expand each category.
Our emphasis on extensibility further enables contributions of new attacks and defenses.

\method also includes components for datasets, models, and utilities that enable end-to-end evaluation.

\noindent\textbf{Datasets.}
\method provides five widely used datasets spanning vision and tabular domains:
\fmnist~\cite{xiao2017FashionMNISTNovelImage}, \cifar, \census~\cite{kohavi1996CensusIncome}, \lfw~\cite{melis2019ExploitingUnintendedFeature}, and \celeba~\cite{liu2015DeepLearningFace}.
These datasets cover small grayscale images, natural images, demographic records, and large-scale face attributes,
supporting evaluation across security, privacy, and fairness.
Section~\ref{sec:extensibility} extends this set to text with \sstwo~\cite{socher2013RecursiveDeepModels}.

\noindent\textbf{Model configurations.}
\method includes commonly used architectures for quick testing, such as VGG~\cite{liu2015VeryDeepConvolutional} and ResNet~\cite{he2016DeepResidualLearning}.
Section~\ref{sec:extensibility} adds an LLM victim, a LoRA-adapted \llama~\cite{llamateammetaai2024Llama3Herd}.

\subsection{Consistent}
Each risk in \method follows a uniform interface: attacks and defenses take a PyTorch \texttt{nn.Module} and hyperparameters as input (except poisoning attacks, which operate on datasets), and their outputs serve as inputs to the corresponding metrics.
This consistent design allows \method to integrate seamlessly into any PyTorch pipeline.

Furthermore, \method offers helper functions for setting up an end-to-end pipeline, including data loading, model initialization, training, and evaluation.
Each helper function, along with every attack and defense, provides default values that yield reasonable results when used with a VGG architecture on \celeba or \cifar.

Attacks, defenses, and metrics follow a consistent layout:
\begingroup
\setlength{\topsep}{0pt}
\setlength{\partopsep}{0pt}
\begin{quote}
  \texttt{amulet.<risk>.attacks.<attack>} \\
  \texttt{amulet.<risk>.defenses.<defense>} \\
  \texttt{amulet.<risk>.metrics.<metric>}
\end{quote}
\endgroup

\input{figures/fig_code}

Figure~\ref{fig:code} illustrates example usage, assuming standard variables common to most ML pipelines.
Each algorithm is instantiated as an object and then executed via a function call.
More detailed examples are available on GitHub.\footnote{\url{https://github.com/ssg-research/amulet/tree/main/examples}}
% More detailed examples are available on GitHub.\footnote{Link shared in supplemental material.}

\method also provides an \texttt{AmuletDataset} class.
Converting a dataset into this format makes it compatible with any module in \method.
All modules support image-based (two-dimensional), tabular (one-dimensional), and text datasets.
Section~\ref{sec:extensibility} describes how we added the text modality and what it cost.

\begin{takeaway}
One interface exposes every risk's modules, so a defense for one risk and an attack for another compose into a single measurable pipeline.
\end{takeaway}
